% GENERATED FILE. Edit canonical lessons, then run npm run generate:books.
% canonical-source-hash: fnv1a64:3f99d05006cee64f
% canonical-lessons: UR-C394-garha, UR-C394-draivar, UR-C394-dhakka, UR-C394-uchhalna, UR-C394-viran

\chapter{Potholes and Jolts}
\label{ch:ur-potholes-and-jolts}
\hlchaptermodality{394}

\begin{chapteropening}
\textbf{By the end of this chapter:} \emph{I can describe a bumpy road: potholes, jolts and an empty stretch at night.}

Complete the last lesson of chapter 394: I can describe a bumpy road: potholes, jolts and an empty stretch at night.
\end{chapteropening}

\section[\texorpdfstring{\ur{گڑھا} (gaṛhā)}{گڑھا (gaṛhā)}]{\ur{گڑھا} (gaṛhā) — a pothole; a pit}

\label{lesson:UR-C394-garha}



\begin{quote}
Before the new one: say the Urdu for to edge forward, then the Urdu for to have (something) stopped.
\end{quote}

\subsection*{You'll want to know: \ur{گڑھا}}
\textbf{\ur{گڑھا}} — \emph{gaṛhā} — \textquotedblleft{}a pothole; a pit\textquotedblright{}.

\begin{quote}
\ur{اس} \ur{سڑک} \ur{میں} \ur{بڑا} \ur{گڑھا} \ur{ہے۔}
\end{quote}

\begin{quote}
\emph{is saṛak meṅ baṛā gaṛhā hai.}
\end{quote}

\begin{quote}
There is a big pothole in this road.
\end{quote}

Masculine.

\begin{cousinweb}
\textbf{\ur{گڑھا}} (\emph{gaṛhā}) probably belongs with the family CDIAL files under \emph{gaḍḍa}, \textquotedblleft{}a hole, a pit\textquotedblright{}, a form it reconstructs: CDIAL lists the Hindi cousin \emph{gāṛā}, \textquotedblleft{}a hole\textquotedblright{}, there (CDIAL 3981), but not this form with \emph{ṛh}.
\end{cousinweb}

\begin{grammarlens}[title={What you've built}]
One more word for getting around.
\end{grammarlens}

\subsection*{Your turn}
\begin{itemize}
\raggedright
  \item \emph{Say it:} \emph{gaṛhā}
  \item \emph{Say it:} \emph{gaṛhā}, once more
  \item \emph{Say it:} \emph{is saṛak meṅ baṛā gaṛhā hai.} — There is a big pothole in this road.
  \item \emph{Recall:} say the Urdu for to edge forward, then the Urdu for to have (something) stopped, then say \emph{gaṛhā} again
\end{itemize}

\begin{tcolorbox}[breakable,colback=teal!4,colframe=teal!35!black,title={Before you move on}]
What does \ur{گڑھا} mean? (\textquotedblleft{}A pothole; a pit\textquotedblright{}.) Say it once more.
\end{tcolorbox}
\section[\texorpdfstring{\ur{ڈرائیور} (ḍrāivar)}{ڈرائیور (ḍrāivar)}]{\ur{ڈرائیور} (ḍrāivar) — a driver}

\label{lesson:UR-C394-draivar}



\begin{quote}
Before the new one: say the Urdu for to have (something) stopped, then the Urdu for a pothole.
\end{quote}

\subsection*{You'll want to know: \ur{ڈرائیور}}
\textbf{\ur{ڈرائیور}} — \emph{ḍrāivar} — \textquotedblleft{}a driver\textquotedblright{}.

\begin{quote}
\ur{ڈرائیور} \ur{نے} \ur{گڑھا} \ur{نہیں} \ur{دیکھا۔}
\end{quote}

\begin{quote}
\emph{ḍrāivar ne gaṛhā nahīṅ dekhā.}
\end{quote}

\begin{quote}
The driver didn't see the pothole.
\end{quote}

Masculine or feminine, by the person meant.

\begin{cousinweb}
\textbf{\ur{ڈرائیور}} (\emph{ḍrāivar}) is borrowed from English \emph{driver}.
\end{cousinweb}

\begin{grammarlens}[title={What you've built}]
One more word for getting around.
\end{grammarlens}

\subsection*{Your turn}
\begin{itemize}
\raggedright
  \item \emph{Say it:} \emph{ḍrāivar}
  \item \emph{Say it:} \emph{ḍrāivar}, once more
  \item \emph{Say it:} \emph{ḍrāivar ne gaṛhā nahīṅ dekhā.} — The driver didn't see the pothole.
  \item \emph{Recall:} say the Urdu for to have (something) stopped, then the Urdu for a pothole, then say \emph{ḍrāivar} again
\end{itemize}

\begin{tcolorbox}[breakable,colback=teal!4,colframe=teal!35!black,title={Before you move on}]
What does \ur{ڈرائیور} mean? (\textquotedblleft{}A driver\textquotedblright{}.) Say it once more.
\end{tcolorbox}
\section[\texorpdfstring{\ur{دھکا} (dhakkā)}{دھکا (dhakkā)}]{\ur{دھکا} (dhakkā) — a shove; a jolt}

\label{lesson:UR-C394-dhakka}



\begin{quote}
Before the new one: say the Urdu for a pothole, then the Urdu for a driver.
\end{quote}

\subsection*{You'll want to know: \ur{دھکا}}
\textbf{\ur{دھکا}} — \emph{dhakkā} — \textquotedblleft{}a shove; a jolt\textquotedblright{}.

\begin{quote}
\ur{گڑھے} \ur{کی} \ur{وجہ} \ur{سے} \ur{بس} \ur{کو} \ur{دھکا} \ur{لگا۔}
\end{quote}

\begin{quote}
\emph{gaṛhe kī vajah se bas ko dhakkā lagā.}
\end{quote}

\begin{quote}
The bus gave a jolt because of the pothole.
\end{quote}

Masculine.

\begin{cousinweb}
\textbf{\ur{دھکا}} (\emph{dhakkā}) is inherited from \emph{dhakk}, \textquotedblleft{}push, strike\textquotedblright{}, a form CDIAL reconstructs (CDIAL 6701); \textbf{\ur{دھکیلنا}}, \textquotedblleft{}to push\textquotedblright{}, comes from the same source.
\end{cousinweb}

\begin{grammarlens}[title={What you've built}]
One more word for getting around.
\end{grammarlens}

\subsection*{Your turn}
\begin{itemize}
\raggedright
  \item \emph{Say it:} \emph{dhakkā}
  \item \emph{Say it:} \emph{dhakkā}, once more
  \item \emph{Say it:} \emph{gaṛhe kī vajah se bas ko dhakkā lagā.} — The bus gave a jolt because of the pothole.
  \item \emph{Recall:} say the Urdu for a pothole, then the Urdu for a driver, then say \emph{dhakkā} again
\end{itemize}

\begin{tcolorbox}[breakable,colback=teal!4,colframe=teal!35!black,title={Before you move on}]
What does \ur{دھکا} mean? (\textquotedblleft{}A shove; a jolt\textquotedblright{}.) Say it once more.
\end{tcolorbox}
\section[\texorpdfstring{\ur{اچھلنا} (uchhalnā)}{اچھلنا (uchhalnā)}]{\ur{اچھلنا} (uchhalnā) — to bounce; to jump up}

\label{lesson:UR-C394-uchhalna}



\begin{quote}
Before the new one: say the Urdu for a driver, then the Urdu for a shove.
\end{quote}

\subsection*{You'll want to know: \ur{اچھلنا}}
\textbf{\ur{اچھلنا}} — \emph{uchhalnā} — \textquotedblleft{}to bounce; to jump up\textquotedblright{}.

\begin{quote}
\ur{دھکے} \ur{سے} \ur{سب} \ur{مسافر} \ur{سیٹوں} \ur{پر} \ur{اچھلے۔}
\end{quote}

\begin{quote}
\emph{dhakke se sab musāfir sīṭoṅ par uchhle.}
\end{quote}

\begin{quote}
With the jolt, all the passengers bounced in their seats.
\end{quote}

Past (he …): \ur{وہ} \ur{سیٹ} \ur{پر} \ur{اچھلا}, \emph{voh sīṭ par uchhlā}.

\begin{cousinweb}
\textbf{\ur{اچھلنا}} (\emph{uchhalnā}) is inherited: CDIAL lists the Hindi cousin \emph{uchalnā}, \textquotedblleft{}to leap\textquotedblright{}, under \emph{ut-śalati}, \textquotedblleft{}springs up\textquotedblright{}, a form it reconstructs (CDIAL 1843).
\end{cousinweb}

\begin{grammarlens}[title={What you've built}]
One more word for getting around.
\end{grammarlens}

\subsection*{Your turn}
\begin{itemize}
\raggedright
  \item \emph{Say it:} \emph{uchhalnā}
  \item \emph{Say it:} \emph{uchhalnā}, once more
  \item \emph{Say it:} \emph{dhakke se sab musāfir sīṭoṅ par uchhle.} — With the jolt, all the passengers bounced in their seats.
  \item \emph{Recall:} say the Urdu for a driver, then the Urdu for a shove, then say \emph{uchhalnā} again
\end{itemize}

\begin{tcolorbox}[breakable,colback=teal!4,colframe=teal!35!black,title={Before you move on}]
What does \ur{اچھلنا} mean? (\textquotedblleft{}To bounce; to jump up\textquotedblright{}.) Say it once more.
\end{tcolorbox}
\section[\texorpdfstring{\ur{ویران} (vīrān)}{ویران (vīrān)}]{\ur{ویران} (vīrān) — deserted; desolate}

\label{lesson:UR-C394-viran}



\begin{quote}
Before the new one: say the Urdu for a shove, then the Urdu for to bounce.
\end{quote}

\subsection*{You'll want to know: \ur{ویران}}
\textbf{\ur{ویران}} — \emph{vīrān} — \textquotedblleft{}deserted; desolate\textquotedblright{}.

\begin{quote}
\ur{رات} \ur{کو} \ur{یہ} \ur{سڑک} \ur{ویران} \ur{ہوتی} \ur{ہے۔}
\end{quote}

\begin{quote}
\emph{rāt ko yih saṛak vīrān hotī hai.}
\end{quote}

\begin{quote}
At night this road is deserted.
\end{quote}

\begin{cousinweb}
is usually traced to Persian \emph{vīrān}, \textquotedblleft{}ruined, laid waste\textquotedblright{}; in Urdu an empty road or an abandoned house is \textbf{\ur{ویران}}.
\end{cousinweb}

\begin{grammarlens}[title={What you've built}]
One more word for getting around.
\end{grammarlens}

\subsection*{Your turn}
\begin{itemize}
\raggedright
  \item \emph{Say it:} \emph{vīrān}
  \item \emph{Say it:} \emph{vīrān}, once more
  \item \emph{Say it:} \emph{rāt ko yih saṛak vīrān hotī hai.} — At night this road is deserted.
  \item \emph{Recall:} say the Urdu for a shove, then the Urdu for to bounce, then say \emph{vīrān} again
\end{itemize}

\begin{tcolorbox}[breakable,colback=teal!4,colframe=teal!35!black,title={Before you move on}]
What does \ur{ویران} mean? (\textquotedblleft{}Deserted; desolate\textquotedblright{}.) Say it once more.
\end{tcolorbox}
